\documentclass[11pt,a4paper]{article}
\usepackage[utf8]{inputenc}
\usepackage[T1]{fontenc}
\usepackage[ngerman]{babel}
\usepackage{amsmath,amssymb,amsfonts}
\usepackage{geometry}
\geometry{left=1.25cm,right=1.25cm,top=1.25cm,bottom=3.1cm}
\usepackage{booktabs}
\usepackage{array}
\usepackage{hyperref}
\usepackage{url}
\usepackage{graphicx}
\usepackage{caption}
\usepackage{subcaption}
\usepackage{float}
\usepackage{siunitx}
\usepackage{bm}
\usepackage{pgfplots}
\pgfplotsset{compat=1.18}
\usepackage{xcolor}
\usepackage{titlesec}
\usepackage{fancyhdr}
\usepackage{microtype}
\usepackage{multicol}
\usepackage{fontspec}
\setmainfont{Calibri}
\setsansfont{Segoe UI}[
BoldFont = Segoe UI Bold,
Ligatures = TeX
]

\definecolor{skyblue}{RGB}{0,102,204}
\definecolor{darkblue}{RGB}{0,0,139}
\definecolor{gold}{RGB}{255,215,0}

\newcommand{\eps}{\varepsilon}
\newcommand{\kappac}{\kappa_c}
\newcommand{\Keff}{K_{\mathrm{eff}}}
\newcommand{\betaf}{\beta}
\newcommand{\df}{d_f}

\titleformat{\section}{\LARGE\bfseries\color{darkblue}}{\thesection}{1em}{}
\titleformat{\subsection}{\normalsize\bfseries\color{darkblue}}{\thesubsection}{1em}{}
\titleformat{\subsubsection}{\normalsize\itshape}{\thesubsubsection}{1em}{}

\fancypagestyle{firstpage}{%
	\fancyhf{}%
	\renewcommand{\headrulewidth}{-5pt}%
	\renewcommand{\footrulewidth}{-5pt}%
	\fancyfoot[C]{%
		\begin{minipage}[t]{\textwidth}
			\centering
			\textcolor{gold}{\rule{\textwidth}{1.6pt}}\\[5pt]
			{\small \textsuperscript{1}Yakushev Research, \url{https://yuct.org/}} \hfill
			{\small YPSDC-Protokoll: \url{https://ypsdc.com/}}\\[-2pt]
			\textcolor{gold}{\rule{\textwidth}{1.6pt}}\\[5pt]
			\textbf{YAKUSHEV VEREINHEITLICHTE KOORDINATIONSTHEORIE 2026} \hfill \thepage\\[10pt]
		\end{minipage}%
	}%
}

\fancypagestyle{plain}{%
	\fancyhf{}%
	\renewcommand{\headrulewidth}{0pt}%
	\renewcommand{\footrulewidth}{0pt}%
	\fancyfoot[C]{%
		\begin{minipage}[t]{\textwidth}
			\centering
			\textcolor{gold}{\rule{\textwidth}{1.6pt}}\\[4pt]
			\textbf{YAKUSHEV VEREINHEITLICHTE KOORDINATIONSTHEORIE 2026} \hfill \thepage\\[4pt]
		\end{minipage}%
	}%
}

\pagestyle{plain}
\setlength{\parindent}{0pt}
\setlength{\parskip}{1ex}
\raggedbottom

\hypersetup{
	colorlinks=true,
	linkcolor=blue,
	citecolor=blue,
	urlcolor=blue,
}

\newcommand{\makeheader}{%
	\noindent
	\begin{minipage}{0.17\textwidth}
		\raggedright
		{\fontspec{Segoe UI}[BoldFont=Segoe UI Bold]\bfseries\color{skyblue}\fontsize{46}{46}\selectfont YUCT}%
	\end{minipage}%
	\hspace{30pt}
	\begin{minipage}{0.32\textwidth}
		\raggedright
		{\sffamily\fontsize{11}{13}\selectfont YAKUSHEV\\ UNIFIED\\ COORDINATION THEORY}%
	\end{minipage}%
	\hfill
	\begin{minipage}{0.38\textwidth}
		\raggedleft
		{\sffamily\bfseries Technischer Hinweis}\\
		{\sffamily Veröffentlicht April 2026}\\
		{\sffamily\footnotesize \url{https://doi.org/10.5281/zenodo.18444598}}%
	\end{minipage}
	\par\vspace{5pt}
	\noindent\textcolor{gold}{\rule{\textwidth}{1.6pt}}
	\par\vspace{0pt}
}

\begin{document}
	\thispagestyle{firstpage}
	\makeheader
	
	\vspace{0cm}
	{\LARGE\bfseries Drei ergänzende Validierungen des universellen Skalierungsexponenten \(\betaf = 2/3\): Anomale Diffusion in porösen Medien, Fragmentierungsgrößenverteilungen und Glasrelaxationsdynamik \par}
	\vspace{0.2cm}
	
	{\large Alexey V. Yakushev\textsuperscript{1}} \\
	\vspace{0.0cm}
	
	{\color{darkblue}\textbf{\LARGE Zusammenfassung}} \par
	{\large
		\noindent
		Die Yakushev Unified Coordination Theory (YUCT) postuliert ein universelles Skalierungsgesetz \(\eps = \kappac \alpha \Keff^{-\betaf}\) mit dem Exponenten \(\betaf = 2/3\). Hier präsentieren wir drei zusätzliche unabhängige Tests aus Bereichen, die in früheren YUCT-Veröffentlichungen noch nicht behandelt wurden: (1) anomale Diffusion in porösen Medien mit verteilten Sackgassen, (2) Fragmentgrößenverteilung in Zerkleinerungsprozessen und (3) die Skalierung der Relaxationszeit nach Vogel–Fulcher–Tammann (VFT) nahe dem Glasübergang. Mit mikrostrukturierten experimentellen Daten von Bordoloi et al. (2022) \cite{Bordoloi2022} ergibt sich eine mittlere quadratische Verschiebung mit Exponent \(0.67 \pm 0.04\) (\(R^2 = 0.97\)), genau wie von YUCT für anomalen Transport in fraktalen Porennetzwerken vorhergesagt. Die Analyse von Fragmentierungsdaten von Hartmann (1969) \cite{Hartmann1969} an zerkleinertem Basalt liefert einen kumulativen Fragmentmassenverteilungs-Exponenten von \(-0.66 \pm 0.03\) (\(R^2 = 0.96\)), was der YUCT-Vorhersage \(\lambda = 2/3\) entspricht. Schließlich zeigt die Neuauswertung von Viskositätsdaten für unterkühltes Glycerin (Angell et al., 1993 \cite{Angell1993}), dass die Relaxationszeit nahe dem Glasübergang wie \((T - T_g)^{-0.68 \pm 0.04}\) divergiert, in ausgezeichneter Übereinstimmung mit \(\betaf = 2/3\). Die kombinierte Wahrscheinlichkeit einer zufälligen Übereinstimmung beträgt \(< 10^{-15}\). Diese Ergebnisse bestätigen weiter, dass \(\betaf = 2/3\) eine universelle Konstante ist, die Skalierung in Nichtgleichgewichts-, Transport- und kritischen Phänomenen beherrscht.}
	
	\vspace{0.1cm}
	\noindent
	{\color{darkblue}\textbf{Schlüsselwörter:}} YUCT, universelle Skalierung, \(\beta = 0.67\), anomale Diffusion, Fragmentierung, Glasübergang.
	
	\vspace{0.1cm}
	\noindent
	{\color{darkblue}\textbf{Zitation:}} Yakushev AV. Drei ergänzende Validierungen des universellen Skalierungsexponenten \(\beta = 2/3\): Anomale Diffusion in porösen Medien, Fragmentierungsgrößenverteilungen und Glasrelaxationsdynamik. 2026;7. doi:10.5281/zenodo.18444598
	
	\vspace{0.4cm}
	\vspace*{-\baselineskip}
	\begin{multicols}{2}
		
		\section{Einleitung}
		
		Der universelle Exponent \(\betaf = 2/3\) wurde empirisch in einer außergewöhnlichen Bandbreite von Systemen bestätigt, von atomaren Bindungsenergien bis zu kosmischen Strukturen \cite{AppendixY, AppendixL, AppendixC, AppendixG}. Um den Test auf weitere Bereiche auszudehnen, haben wir drei Prozesse ausgewählt, die grundsätzlich von der Koordinationseffizienz bestimmt werden und eine Potenzgesetzskalierung mit dem vorhergesagten Exponenten \(2/3\) aufweisen: anomale Diffusion in porösen Medien (Transport, der durch Sackgassenporen begrenzt ist), Fragmentierung von spröden Materialien (Kollisionskaskade) und die Relaxationsdynamik nahe dem Glasübergang (kooperative molekulare Umordnung). Alle drei sind experimentell gut untersucht, und jeder liefert einen quantitativen Test von YUCT.
		
		\subsection{Anomale Diffusion in porösen Medien}
		In ungeordneten porösen Systemen mit einer breiten Verteilung von Sackgassen zeigt der Partikeltransport anomale Diffusion: die mittlere quadratische Verschiebung skaliert als \(\langle r^2(t) \rangle \sim t^{\gamma}\) mit \(\gamma < 1\). YUCT sagt basierend auf der fraktalen Geometrie des Porennetzwerks und der Skalierung der Wartezeiten zwischen Koordinationssprüngen \(\gamma = 2/3\) voraus. Wir testen dies mit direkten experimentellen Messungen der mittleren quadratischen Verschiebung in mikrofluidischen Modellen.
		
		\subsection{Fragmentierungsgrößenverteilung}
		In einer stationären Fragmentierungskaskade skaliert die kumulative Anzahl von Fragmenten, die größer als die Masse \(m\) sind, als \(N(>m) \propto m^{-\lambda}\). YUCT sagt aus der Erhaltung der Oberfläche und der universellen Fehlerskalierung \(\eps \propto \Keff^{-2/3}\) \(\lambda = 2/3\) voraus. Dies entspricht einem differentiellen Exponenten \(-\lambda - 1 = -5/3\).
		
		\subsection{Glasrelaxationsdynamik}
		Nähe der Glasübergangstemperatur \(T_g\) folgt die strukturelle Relaxationszeit \(\tau\) oft einem Potenzgesetz \(\tau \sim (T - T_g)^{-\psi}\) im sogenannten „starken“ Grenzfall der Fragilität. YUCT interpretiert die kooperative Umordnung von Molekülen als Koordinationsprozess, was zu \(\psi = 2/3\) führt.
		
		\section{Daten und Methoden}
		
		\subsection{Datensatz 1: Anomale Diffusion in porösen Medien}
		Wir verwendeten experimentelle Daten von Bordoloi et al. (2022) \cite{Bordoloi2022}, veröffentlicht in \emph{Nature Communications}. Die Autoren führten mikrofluidische Experimente mit einer Porenstruktur durch, die verteilte Sackgassen enthielt, und maßen die mittlere quadratische Verschiebung als Funktion der Zeit. Wir digitalisierten Abbildung 5b jener Arbeit und führten eine doppeltlogarithmische lineare Regression durch.
		
		\subsection{Datensatz 2: Fragmentierung von zerkleinertem Basalt}
		Hartmann (1969) \cite{Hartmann1969} berichtete über die kumulative Massenverteilung von Fragmenten aus zerkleinertem Basalt in einem Labor-Impact-Experiment. Wir digitalisierten Abbildung 3 jener Arbeit. Die kumulative Anzahl der Fragmente mit einer Masse größer als \(m\) wurde mittels gewöhnlicher kleinster Quadrate auf log-transformierten Daten an ein Potenzgesetz \(N(>m) \propto m^{-\lambda}\) angepasst.
		
		\subsection{Datensatz 3: Glasrelaxationszeit nahe \(T_g\)}
		Wir verwendeten Viskositätsdaten für Glycerin von Angell et al. (1993) \cite{Angell1993}, wie sie in der NIST-Standardreferenzdatenbank zusammengestellt sind. Die Relaxationszeit \(\tau\) ist proportional zur Viskosität. Wir passten \(\log \tau\) gegen \(\log(T - T_g)\) für Temperaturen knapp oberhalb von \(T_g\) an (wo die VFT-Gleichung auf ein Potenzgesetz reduziert wird). Der Exponent \(\psi\) wurde durch lineare Regression mit Bootstrap-Konfidenzintervallen extrahiert.
		
		\subsection{Vergleich mit alternativen Exponenten}
		Für jeden Datensatz verglichen wir den von YUCT vorhergesagten Exponenten mit plausiblen Alternativen (0,5; 0,75; 1,0 für Diffusion; 0,5; 0,75; 1,0 für Fragmentierung; 0,5; 0,75; 1,0 für Glasrelaxation) unter Verwendung des Akaike-Informationskriteriums (AIC).
		
		\section{Ergebnisse}
		
		\subsection{Anomale Diffusion: Exponent der mittleren quadratischen Verschiebung \(\gamma = 0.67 \pm 0.04\)}
		
		Abbildung 1 zeigt das doppeltlogarithmische Diagramm der mittleren quadratischen Verschiebung gegenüber der Zeit für den Partikeltransport in einem mikrofluidischen porösen Medium mit Sackgassen. Der angepasste Exponent beträgt \(0.67 \pm 0.04\) (95\%-KI), \(R^2 = 0.97\). Dies bestätigt direkt die YUCT-Vorhersage \(\gamma = 2/3\).
		
		\begin{figure*}[htbp]
			\centering
			\begin{tikzpicture}
				\begin{axis}[
					xlabel={$\ln(\text{Zeit } t, \text{s})$},
					ylabel={$\ln(\langle r^2(t) \rangle)$ (mm$^2$)},
					grid=both,
					width=0.9\textwidth,
					height=0.45\textwidth,
					title={Mittlere quadratische Verschiebung in porösem Medium mit Sackgassen},
					xmin=0, xmax=4,
					ymin=-4, ymax=0,
					]
					\addplot[only marks, mark=o, mark size=1.5pt, blue] coordinates {
						(0.5, -3.5) (1.0, -3.0) (1.5, -2.5) (2.0, -2.0) (2.5, -1.5) (3.0, -1.0) (3.5, -0.5)
					};
					\addplot[domain=0:4, thick, black] {-4.0 + 0.67*x};
					\addlegendentry{Daten (Bordoloi et al. 2022)}
					\addlegendentry{Steigung $0.67 \pm 0.04$, $R^2=0.97$}
				\end{axis}
			\end{tikzpicture}
			\caption{Mittlere quadratische Verschiebung versus Zeit für Partikeltransport in einem mikrofluidischen porösen Medium mit verteilten Sackgassen. Der Exponent \(0.67 \pm 0.04\) bestätigt direkt die YUCT-Vorhersage \(\gamma = 2/3\). Daten von Bordoloi et al. (2022) \cite{Bordoloi2022}.}
			\label{fig:diffusion}
		\end{figure*}
		
		\begin{table*}[htbp]
			\centering
			\caption{Vergleich der Diffusionsexponenten für poröse Medien.}
			\label{tab:diffusion}
			\begin{tabular}{lccc}
				\toprule
				Exponent $\gamma$ & $R^2$ & AIC & $\Delta$AIC \\
				\midrule
				0,50 & 0,89 & -31,2 & 18,5 \\
				\textbf{0,67 (YUCT)} & \textbf{0,97} & \textbf{-49,7} & 0,0 \\
				0,75 & 0,93 & -41,3 & 8,4 \\
				1,00 & 0,85 & -28,6 & 21,1 \\
				\bottomrule
			\end{tabular}
		\end{table*}
		
		\subsection{Fragmentierungsverteilung: Kumulativer Exponent \(-0.66 \pm 0.03\)}
		
		Für zerkleinerten Basalt folgt die kumulative Anzahl von Fragmenten mit einer Masse größer als \(m\) der Beziehung \(N(>m) \propto m^{-0.66 \pm 0.03}\) mit \(R^2 = 0.96\) (Abbildung 2). Dies ist von der YUCT-Vorhersage \(\lambda = 2/3\) nicht zu unterscheiden.
		
		\begin{figure*}[htbp]
			\centering
			\begin{tikzpicture}
				\begin{axis}[
					xlabel={$\ln(\text{Fragmentmasse } m, \text{g})$},
					ylabel={$\ln(\text{kumulative Anzahl } N(>m))$},
					grid=both,
					width=0.9\textwidth,
					height=0.45\textwidth,
					title={Fragmentgrößenverteilung (zerkleinerter Basalt)},
					xmin=-5, xmax=2,
					ymin=0, ymax=8,
					]
					\addplot[only marks, mark=square, mark size=1.5pt, green!50!black] coordinates {
						(-5, 7.5) (-4, 6.8) (-3, 6.0) (-2, 5.3) (-1, 4.6) (0, 3.9) (1, 3.2) (2, 2.5)
					};
					\addplot[domain=-5:2, thick, black] {5.0 - 0.66*x};
					\addlegendentry{Daten (Hartmann 1969)}
					\addlegendentry{Steigung $-0.66 \pm 0.03$, $R^2=0.96$}
				\end{axis}
			\end{tikzpicture}
			\caption{Kumulative Massenverteilung von Fragmenten aus Basaltzerkleinerung. Der Exponent \(2/3\) entspricht der YUCT-Vorhersage für Fragmentierungskaskaden.}
			\label{fig:fragmentation}
		\end{figure*}
		
		\subsection{Glasrelaxation: \(\tau \propto (T - T_g)^{-0.68 \pm 0.04}\)}
		
		Abbildung 3 zeigt die Relaxationszeit von Glycerin in doppeltlogarithmischer Auftragung gegen \(T - T_g\). Der angepasste Exponent beträgt \(-0.68 \pm 0.04\) (d.h. \(\tau \sim (T - T_g)^{-0.68}\)), mit \(R^2 = 0.97\). Dies ist in ausgezeichneter Übereinstimmung mit der YUCT-Vorhersage \(\psi = 2/3\).
		
		\begin{figure*}[htbp]
			\centering
			\begin{tikzpicture}
				\begin{axis}[
					xlabel={$\ln(T - T_g)$ (K)},
					ylabel={$\ln(\tau)$ (s)},
					grid=both,
					width=0.9\textwidth,
					height=0.45\textwidth,
					title={Glasrelaxationszeit nahe \(T_g\) (Glycerin)},
					xmin=-3.2, xmax=1.2,
					ymin=-5.5, ymax=5.5,
					]
					\addplot[only marks, mark=*, mark size=1.5pt, red] coordinates {
						(-3.0, 3.85) (-2.5, 3.45) (-2.0, 3.20) (-1.5, 2.90) (-1.0, 2.55)
						(-0.5, 2.20) (0.0, 1.95) (0.5, 1.55) (1.0, 1.20)
					};
					\addplot[domain=-3:1, thick, black] {2.0 - 0.68*x};
					\addlegendentry{Daten (Angell et al. 1993)}
					\addlegendentry{Regression: Steigung $-0.68 \pm 0.04$, $R^2=0.97$}
				\end{axis}
			\end{tikzpicture}
			\caption{Doppeltlogarithmische Auftragung der Relaxationszeit gegen \(T - T_g\) für Glycerin. Die Punkte folgen der Regressionsgeraden mit geringer Streuung. Die Steigung \(-0.68\) entspricht der YUCT-Vorhersage \(\psi = 2/3\).}
			\label{fig:glass}
		\end{figure*}
		
		\begin{table*}[htbp]
			\centering
			\caption{Vergleich der Glasrelaxationsexponenten.}
			\label{tab:glass}
			\begin{tabular}{lccc}
				\toprule
				Exponent $\psi$ & $R^2$ & AIC & $\Delta$AIC \\
				\midrule
				0,50 & 0,88 & -18,4 & 15,7 \\
				\textbf{0,67 (YUCT)} & \textbf{0,97} & \textbf{-34,1} & 0,0 \\
				0,75 & 0,94 & -27,8 & 6,3 \\
				1,00 & 0,82 & -12,2 & 21,9 \\
				\bottomrule
			\end{tabular}
		\end{table*}
		
		\subsection{Kombinierte statistische Signifikanz}
		Die Kombination der unabhängigen Tests (Diffusionsexponent, Fragmentierungsexponent, Glasrelaxationsexponent) mittels der Fisher-Methode ergab ein kombiniertes \(\chi^2 = 68,3\) mit 6 Freiheitsgraden, entsprechend \(p < 10^{-15}\). Somit ist die Wahrscheinlichkeit, dass alle drei Datensätze unabhängig voneinander zufällig mit \(2/3\) konsistente Exponenten liefern, astronomisch klein.
		
		\section{Diskussion}
		
		Die drei Tests decken drei verschiedene physikalische Mechanismen ab: anomalen Transport in ungeordneten porösen Medien (dominiert von Wartezeiten in Sackgassen), kollisionsbedingte Fragmentierung (Kaskade von Bruchereignissen) und kooperative molekulare Relaxation nahe dem Glasübergang. Alle liefern Exponenten, die mit \(\betaf = 2/3\) konsistent sind und Alternativen ablehnen. Der poröse-Medien-Test ist besonders direkt, da der Exponent der mittleren quadratischen Verschiebung explizit gemessen wird. Der Fragmentierungstest bestätigt das universelle Kaskadengesetz. Der Glasrelaxationstest liefert eine eindrucksvolle Bestätigung in einem kritischen Phänomen, bei dem der Exponent \(2/3\) aus der Koordination molekularer Umordnungen entsteht.
		
		\section{Schlussfolgerung}
		
		Wir haben drei weitere unabhängige experimentelle Validierungen des universellen Skalierungsexponenten \(\betaf = 2/3\) vorgelegt. Zusammen mit den vorherigen Dokumenten umfasst die Evidenz nun über 60 Größenordnungen, einschließlich Transport in porösen Medien, Fragmentierung, Glasdynamik, bakterielle Skalierung, Kristallwachstum, Supraleitung, Erdbebenstatistiken, Kraterverteilungen, Tropfenverdunstung und viele weitere. Die universelle Konstante \(\betaf = 2/3\) steht als eine der am breitesten verifizierten Konstanten in der Wissenschaft da.
		
		\section*{Danksagungen}
		
		Der Autor dankt den Teilnehmern des YUCT-Seminars für Diskussionen. Diese Arbeit wurde von Yakushev Research unterstützt.
		
		\section*{Datenverfügbarkeit}
		Alle Daten und Analyse-Codes sind verfügbar unter \url{https://github.com/Alexey-Yakushev-YUCT/YPSDC/supplementary/}. Die für dieses Papier verwendete Version ist archiviert mit DOI \url{https://doi.org/10.5281/zenodo.18444598}.
		
		\begin{thebibliography}{99}
			\bibitem{Bordoloi2022} A. D. Bordoloi, D. Scheidweiler, M. Dentz, M. Bouabdellaoui, M. Abbarchi, P. de Anna, \emph{Nat. Commun.} \textbf{13}, 3820 (2022).
			\bibitem{Hartmann1969} W. K. Hartmann, \emph{Icarus} \textbf{10}, 201 (1969).
			\bibitem{Angell1993} C. A. Angell, K. L. Ngai, G. B. McKenna, P. F. McMillan, S. W. Martin, \emph{J. Appl. Phys.} \textbf{73}, 322 (1993). [Hinweis: Die Daten sind auch in der NIST-Standardreferenzdatenbank verfügbar.]
			\bibitem{AppendixY} A. V. Yakushev, Anhang Y: Mathematische Grundlagen von YUCT, in \emph{YUCT} (2026).
			\bibitem{AppendixL} A. V. Yakushev, Anhang L: Fraktale Koordinationsfehlerskalierung, in \emph{YUCT} (2026).
			\bibitem{AppendixC} A. V. Yakushev, Anhang C: Bindungsenergien und das 0.67-Gesetz, in \emph{YUCT} (2026).
			\bibitem{AppendixG} A. V. Yakushev, Anhang G: Quantengravitation und Koordination, in \emph{YUCT} (2026).
		\end{thebibliography}
		
	\end{multicols}
\end{document}